% part: intuitionistic-logic
% chapter: propositions-as-types
% section: terms-to-proofs

\documentclass[../../../include/open-logic-section]{subfiles}

\begin{document}

\olfileid{int}{pty}{tp}

\olsection{له ثبوت ترمونو څخه \usetoken{P}{derivation} بېرته اخيستل}

اوس بل لور ته ګورو: يو ثبوت ترم~$N$ بېرته
په \Log{N2i} کښې !!a{proof} ته اړول۔ په عمومي
ډول، دا يوازې د داسې سياق په نسبت شوني دي
چې د ورکړي ثبوت ترم د هر آزاد متغير~$x$ دپاره
د $x : !A$ جوړه ولري۔ خو لومړے داسې بېلګه
اخلو چې آزاد متغيرونه نۀ لري؛ يعنې دا ترم:
\[
\lambd[\typeof{x}{!A \land
  !B}][\pair{\proj{2}{x}}{\proj{1}{x}}]
\]
دا د $\lambd[\typeof{x}{!A \land !B}][N_1]$ په
بڼه دے۔ په \Log{tN2i} کښې يوازينۍ قاعده چې
داسې ترم جوړوي،~\Intro{\lif} ده۔ نو پوهېږو چې
د~$N$ کوډ شوے !!{proof} بايد په~\Intro{\lif}
پاې ته ورسېږي؛ يعنې په دې پاې ته رسېږي:
\[
  \Axiom$x : !A\land !B \fCenter \pair{\proj{2}{x}}{\proj{1}{x}} : 
    !C$
  \RightLabel{\Intro\lif}
  \UnaryInf$\fCenter \lambd[\typeof{x}{!A \land
  !B}][\pair{\proj{2}{x}}{\proj{1}{x}}] : (!A \land !B) \lif !C$
  \DisplayProof
  \]
تر اوسه نۀ پوهېږو چې $!C$ څه دے، خو دا پوهېږو
چې مخکنۍ برخه $!A \land !B$ ده، او د مقدمې په
سياق کښې $x:!A \land !B$ شاملېدل پکار دي۔

اوس د مقدمې اړوند ثبوت ترم له مخه ټاکل شوے
دے: $\pair{\proj{2}{x}}{\proj{1}{x}}$۔ هغه
!!{proof} چې \emph{همدغه} ترم ئې کوډ کوي، بايد
په \Intro{\land} پاې ته ورسېږي۔ لا نۀ پوهېږو
چې $!C$ څه دے، خو څرنګه چې د~\Intro{\land}
پايله ده، بايد د او تړنه وي؛ يعنې
$!C \ident (!C_1 \land !C_2)$۔ بياجوړونه
داسې پرمخ وړو:
\[
  \Axiom$x : !A\land !B \fCenter \proj{2}{x} : !C_1$
  \Axiom$x : !A\land !B \fCenter \proj{1}{x} : !C_2$
  \RightLabel{\Intro\land}
  \BinaryInf$x : !A\land !B \fCenter \pair{\proj{2}{x}}{\proj{1}{x}} : 
    !C_1 \land !C_2$
  \RightLabel{\Intro\lif}
  \UnaryInf$\fCenter \lambd[\typeof{x}{!A \land
  !B}][\pair{\proj{2}{x}}{\proj{1}{x}}] : (!A \land !B) \lif (!C_1 \land !C_2)$
  \DisplayProof
  \]
ثبوت ترم $\pi_2(x)$ ښيي چې پاسنے کيڼ سېکوېنټ
د \Elim{\land}[2] په وسيله ترلاسه شوے؛ يعنې
مقدمه د $!D \land !C_1$ په بڼه ده۔ په ښي‌لاس
اړخ کښې دا ټاکلے شو چې $!C_2$ ته رسوونکے
استنتاج \Elim{\land}[1] دے، او مقدمه ئې د
$!C_2 \land !E$ په بڼه ده۔
\[
  \Axiom$x : !A\land !B \fCenter x: !D\land !C_1$
  \RightLabel{\Elim\land[2]}
  \UnaryInf$x : !A\land !B \fCenter \proj{2}{x} : !C_1$
  \Axiom$x : !A\land !B \fCenter x: !C_2\land !E$
  \RightLabel{\Elim\land[1]}
  \UnaryInf$x : !A\land !B \fCenter \proj{1}{x} : !C_2$
  \RightLabel{\Intro\land}
  \BinaryInf$x : !A\land !B \fCenter \pair{\proj{2}{x}}{\proj{1}{x}} : 
    !C_1 \land !C_2$
  \RightLabel{\Intro\lif}
  \UnaryInf$\fCenter \lambd[\typeof{x}{!A \land
  !B}][\pair{\proj{2}{x}}{\proj{1}{x}}] : (!A \land !B) \lif (!C_1 \land !C_2)$
  \DisplayProof
  \]
څرنګه چې د تر ټولو پاسنيو سېکوېنټونو ثبوت
ترمونه اوس متغيرونه دي، پوهېږو چې اکسيومونه
دي۔ دا $!C_1$، $!C_2$، $!D$ او $!E$ ټاکي:
$!D \ident !A$ او $!C_1 \ident !B$، کنه کيڼ
سېکوېنټ اکسيوم نۀ وي؛ او $!C_2 \ident !A$ او
$!E \ident !B$، کنه ښے سېکوېنټ اکسيوم نۀ وي۔
ثبوت مو په بشپړ ډول بېرته ترلاسه کړ:
\[
  \Axiom$x : !A\land !B \fCenter x: !A\land !B$
  \RightLabel{\Elim\land[2]}
  \UnaryInf$x : !A\land !B \fCenter \proj{2}{x} : !B$
  \Axiom$x : !A\land !B \fCenter x: !A\land !B$
  \RightLabel{\Elim\land[1]}
  \UnaryInf$x : !A\land !B \fCenter \proj{1}{x} : !A$
  \RightLabel{\Intro\land}
  \BinaryInf$x : !A\land !B \fCenter \pair{\proj{2}{x}}{\proj{1}{x}} : 
    !B \land !A$
  \RightLabel{\Intro\lif}
  \UnaryInf$\fCenter \lambd[\typeof{x}{!A \land
  !B}][\pair{\proj{2}{x}}{\proj{1}{x}}] : (!A \land !B) \lif (!B \land !A)$
  \DisplayProof
  \]

\end{document}
